% ====================================================================
% Cadena metodologica del TFM
% ====================================================================
\begin{figure}[H]
\centering
\resizebox{\textwidth}{!}{%
\begin{tikzpicture}[
  font=\sffamily\small,
  stage/.style={draw=VIUDark!70,fill=VIULight,rounded corners=4pt,
    align=center,text width=3.05cm,minimum height=1.25cm,inner sep=6pt},
  key/.style={draw=VIUOrange,fill=VIUOrange!10,rounded corners=4pt,
    align=center,text width=3.05cm,minimum height=1.25cm,inner sep=6pt,
    font=\sffamily\small\bfseries},
  audit/.style={draw=VIUDark!55,fill=VIUWhite,rounded corners=4pt,
    align=center,text width=3.05cm,minimum height=0.9cm,inner sep=5pt,
    font=\sffamily\scriptsize},
  flow/.style={-{Latex[length=2.4mm]},line width=0.9pt,draw=VIUOrange!90!black},
  soft/.style={-{Latex[length=2.0mm]},line width=0.7pt,draw=VIUDark!55}
]
\node[stage] (q) at (0,0) {\textbf{Pregunta}\\coaliciones factibles};
\node[stage] (m) at (3.6,0) {\textbf{Modelo}\\robots\\cargas\\grafo};
\node[stage] (t) at (7.2,0) {\textbf{Teor\'ia}\\utilidad, Smith, cotas};
\node[key]   (e) at (10.8,0) {\textbf{Experimentos}\\SP1--SP8 pareados};
\node[stage] (r) at (14.4,0) {\textbf{Lectura}\\regímenes y límites};

\node[audit] (a1) at (0,-1.95) {fuente bibliografica\\o problema industrial};
\node[audit] (a2) at (3.6,-1.95) {parametros, semillas\\y escenarios};
\node[audit] (a3) at (7.2,-1.95) {resultado algebraico\\o hipótesis};
\node[audit] (a4) at (10.8,-1.95) {CSV, auditorías\\y contraste Holm};
\node[audit] (a5) at (14.4,-1.95) {afirmaci\'on acotada\\sin sobreafirmar};

\foreach \from/\to in {q/m,m/t,t/e,e/r}
  \draw[flow] (\from) -- (\to);
\foreach \from/\to in {q/a1,m/a2,t/a3,e/a4,r/a5}
  \draw[soft] (\from.south) -- (\to.north);

\node[draw=VIUOrange!75,fill=VIUOrange!6,rounded corners=4pt,
  align=center,text width=12.8cm,inner sep=6pt,font=\sffamily\scriptsize,
  below=0.62cm of a3] (rule)
  {Regla de cierre: una afirmación fuerte entra en la memoria solo si apunta a una prueba, una cita, un CSV canónico o una limitación explícita.};
\draw[soft] (a4.south) -- (rule.north);
\end{tikzpicture}}
\caption[Cadena metodológica del TFM]{Cadena metodológica del TFM. La investigación conecta pregunta, modelo, teoría, tratamientos y métricas; la validación no se declara como prueba formal cuando depende de efectos discretos o físicos no cerrados analíticamente.}
\viuownsource
\label{fig:metodologia_investigacion}
\end{figure}
